% ID: SEL-D0031-OA-E03
% ESTADO: cotejada. Segunda pasada visual completa y borrador compilado revisados.
% Fuente: Selectividad/matematicas_aplicadas/titular_junio_6 EXAMEN MAT APLIC C S.pdf
% SHA-256: bebc339cfd8f69c7c4fae160fd165eaef5f6bc93cfd6d15cb9ce289b7029d167
% Página: 1; opción A; ejercicio original 3.

Marta tiene dos trajes rojos, un traje azul y uno blanco. Además, tiene un par de zapatos de color rojo, otro de color azul y dos pares blancos. Si decide aleatoriamente qué ponerse, determine las probabilidades de los siguientes sucesos:
\begin{enumerate}[label=\alph*)]
\item \textbf{(0.8 puntos)} Llevar un traje rojo y unos zapatos blancos.
\item \textbf{(0.9 puntos)} No ir toda vestida de blanco.
\item \textbf{(0.8 puntos)} Calzar zapatos azules o blancos.
\end{enumerate}
